\subsection{The notion of morphism for laws of composition not everywhere defined}

DEFINITION 1. Let E and E' be two sets each with an internal law of composition denoted by \((x,y)\mapsto x*y\) , not necessarily everywhere defined. A mapping \(f:E\to E'\) is a morphism if, for all x,y in E such that \(f(x)*f(y)\) is defined, the composition \(x*y\) is also defined and:

\[
f (x * y) = f (x) * f (y).\tag{1}
\]

More briefly, we may say that formula (1) must hold every time the right hand side has a meaning.

The notion of morphism is distinct from that of representation (Algebra,

Chapter I, § 1, no. 1), where it is demanded that equation (1) hold whenever the left hand side has a meaning. Of course, the two notions coincide for laws of composition everywhere defined.

DEFINITION 2. Let E and E' be two sets each with a family of internal laws of composition \((x,y)\mapsto x*\alpha y,\alpha\in\mathbf{I}\) . A mapping \(f:E\to E'\) is a morphism if it is a morphism for each of the laws of composition \((x,y)\mapsto x*\alpha y\) .

Just as representations, so morphisms satisfy axioms \((\mathrm{MO}_{\mathrm{I}})\) , \((\mathrm{MO}_{\mathrm{II}})\) , \((\mathrm{MO}_{\mathrm{III}})\) of Set Theory, Chapter IV, § 2. If \(f: E \to E'\) is a morphism, \(f(E)\) is a stable subset of \(R'\) .

